\documentclass[11pt]{article}

\usepackage[margin=1in]{geometry}
\usepackage[T1]{fontenc}
\usepackage[utf8]{inputenc}
\usepackage{lmodern}
\usepackage{microtype}
\usepackage{amsmath}
\usepackage{array}
\usepackage{booktabs}
\usepackage{tabularx}
\usepackage{enumitem}
\usepackage[table]{xcolor}
\usepackage{hyperref}
\usepackage{fancyhdr}

\definecolor{ReportBlue}{HTML}{1F4E79}
\definecolor{ReportTeal}{HTML}{1E847F}
\definecolor{ReportGray}{HTML}{F4F6F8}
\definecolor{ReportLine}{HTML}{D7DEE3}

\hypersetup{
  colorlinks=true,
  linkcolor=ReportBlue,
  urlcolor=ReportTeal,
  pdftitle={Skill-JEPA-WM Locked Protocol Progress Report},
  pdfauthor={Codex}
}

\pagestyle{fancy}
\fancyhf{}
\fancyhead[L]{\textcolor{ReportBlue}{Skill-JEPA-WM Locked Progress}}
\fancyhead[R]{2026-04-11}
\fancyfoot[C]{\thepage}
\setlength{\headheight}{14pt}
\setlist[itemize]{leftmargin=1.5em, itemsep=0.35em, topsep=0.4em}

\newcolumntype{Y}{>{\raggedright\arraybackslash}X}

\begin{document}

\begin{center}
{\LARGE\bfseries \textcolor{ReportBlue}{Skill-JEPA-WM Locked Protocol Progress Report}}\\[0.5em]
{\large Push-T Locked Suite Status, public repository snapshot}\\[0.4em]
{\normalsize Scope: reevaluation and incompleteness closure before scaled pilot}
\end{center}

\vspace{0.75em}

\noindent
\fcolorbox{ReportLine}{ReportGray}{
\begin{minipage}{0.97\linewidth}
\textbf{Executive Summary.}
This historical report used the earlier sampled-trajectory goal-state metric.
The later release re-score found \textbf{0\%} standard Push-T coverage success for both planners, so the numbers below are trajectory diagnostics rather than a Push-T task-success claim.
Under that diagnostic, the hierarchical and flat planners tied on sampled goal-state success at \textbf{7\%}, while the hierarchical planner improved mean sampled-state distance from \textbf{355.30} to \textbf{264.43}, reduced average planning latency from \textbf{0.564 s} to \textbf{0.249 s}, and beat flat on \textbf{57\%} of sampled pairs.
In parallel, I fixed the remaining implementation gaps that blocked the scaled pilot: full Push-T pixel decoding, projector-consistent cache generation, exact episode-count cache slicing, a random-skill hierarchical baseline, and locked-suite artifact export.
The 3-seed scaled pilot itself is still incomplete, so no continue/stop decision has been issued for the hierarchy line and no feasibility critic has been added.
\end{minipage}
}

\section*{1. Locked Scope For This Round}

The user request for this round froze the protocol to the corrected report settings:

\begin{itemize}
\item frozen V-JEPA2 cache,
\item \(K = 4\),
\item \texttt{control\_horizon = 32},
\item \texttt{goal\_gap = 24},
\item \texttt{execute\_actions\_per\_plan = 4},
\item cache-consistent online evaluation,
\item deterministic per-episode reseeding.
\end{itemize}

The same request also fixed the experimental boundary:

\begin{itemize}
\item do not add new architecture, loss, or dataset,
\item do not touch DROID,
\item re-evaluate the current best checkpoint first,
\item only add a subgoal feasibility critic if the scaled hierarchy advantage survives.
\end{itemize}

\section*{2. What Was Finished}

\subsection*{2.1 Current Best Checkpoint Reevaluation}

I reran the current best Push-T debug checkpoint on \textbf{100 sampled evaluation pairs} under the locked online evaluator and compared hierarchical vs.\ flat planning from the same checkpoint.
The release re-score later showed that these sampled pairs came from \textbf{one unique episode}, so this artifact is a historical trajectory diagnostic rather than broad held-out task evidence.

\begin{center}
\begin{tabular}{lrrrr}
\toprule
\rowcolor{ReportBlue!10}
Method & Goal-state Diagnostic & Sampled-State Dist. & Planning Latency & Notes \\
\midrule
Hierarchical & 0.07 & 264.43 & 0.249 s & same checkpoint \\
Flat & 0.07 & 355.30 & 0.564 s & same checkpoint \\
\bottomrule
\end{tabular}
\end{center}

Additional comparison summary:

\begin{itemize}
\item hierarchical better-than-flat rate by sampled-state distance: \textbf{0.57},
\item unique episode count after release re-score: \textbf{1},
\item representative rollout GIFs were saved for both planners,
\item per-episode records were exported to CSV.
\end{itemize}

\subsection*{2.2 Incompletenesses Closed Before Scaling}

The main value of this round was not only the reevaluation result.
It was closing the implementation gaps that would have invalidated or blocked the scaled suite.

\begin{itemize}
\item \textbf{Full Push-T read path.} Registered \texttt{hdf5plugin} so the blosc-compressed \texttt{pixels} dataset in the full HDF5 can be read locally.
\item \textbf{Cache consistency.} Reused the frozen projector checkpoint during cache building instead of silently creating a new random projector.
\item \textbf{Scaled cache slicing.} Added exact episode-count trimming via \texttt{max\_episodes} so the cache builder can produce the requested \(12\text{k train} / 1024\text{ val}\) source layout.
\item \textbf{Random baseline.} Added a random-skill hierarchical high-level planner so \texttt{random\_skill\_hier\_10pct} can be evaluated under the same locked suite.
\item \textbf{Artifact export.} Extended the online evaluator to emit per-episode CSV and rollout GIF artifacts for the locked suite.
\item \textbf{Suite entrypoint.} Added a single runner script that prepares configs, launches training and evaluation jobs, and aggregates the scaled suite outputs.
\end{itemize}

\section*{3. Files Added Or Updated}

The main code and artifact paths produced in this round were:

\begin{itemize}
\item \texttt{src/skill\_jepa/analysis/eval\_pusht\_online.py}
\item \texttt{src/skill\_jepa/planning/random\_high\_level.py}
\item \texttt{src/skill\_jepa/planning/\_\_init\_\_.py}
\item \texttt{tools/cache\_vjepa\_features.py}
\item \texttt{tools/run\_skill\_jepa\_pusht\_locked\_suite.py}
\item \texttt{artifacts/pusht\_locked\_suite/evals/current\_best\_checkpoint\_100ep/pusht\_online\_eval.json}
\item \texttt{artifacts/pusht\_locked\_suite/evals/current\_best\_checkpoint\_100ep/pusht\_online\_records.csv}
\item \texttt{artifacts/pusht\_locked\_suite/reports/current\_best\_checkpoint\_comparison.csv}
\end{itemize}

\section*{4. What Is Still Incomplete}

The locked scaled pilot has \textbf{not} been completed yet.
That means the following deliverables are still missing for the scaled suite itself:

\begin{itemize}
\item the full 3-seed training runs on the five requested configurations,
\item aggregate CSV over the scaled pilot,
\item scaled-suite plots,
\item scaled-suite rollout exports,
\item \texttt{decision\_report.md} with a valid continue/stop recommendation,
\item any feasibility-critic implementation or rerun.
\end{itemize}

This stop point was deliberate.
The scaled pilot should not be claimed complete until the actual 3-seed locked runs finish.

\section*{5. Why The Scaled Pilot Was Not Finished In This Round}

I measured the local path before launching the long run.
Using the current local RTX 4070 8GB throughput, the locked suite is large enough that it no longer fits inside a normal interactive turn:

\begin{center}
\begin{tabular}{lr}
\toprule
Estimated Component & Local Time \\
\midrule
Full 13,024-episode cache build & 8.97 h \\
Per-seed training total & 19.93 h \\
3-seed locked suite total & 68.75 h \\
\bottomrule
\end{tabular}
\end{center}

That estimate is large enough to change the execution regime.
It also means the absence of scaled results is a true missing result, not a hidden success.

\section*{6. Interpretation Of The New 100-Episode Result}

The new reevaluation materially weakens the earlier short-run diagnostic.
On the larger 100-episode test:

\begin{itemize}
\item hierarchy keeps a meaningful pose-distance advantage,
\item hierarchy keeps a large latency advantage,
\item hierarchy does \textbf{not} improve standard Push-T coverage success over flat on the current best debug checkpoint.
\end{itemize}

So the current evidence supports this narrower claim:

\begin{quote}
Under the locked debug checkpoint, hierarchy improves sampled-state distance and planning efficiency, but the stronger claim that hierarchy improves standard Push-T coverage success is not supported by the 100-episode reevaluation.
\end{quote}

That is exactly why the scaled pilot remains necessary before any new model component is added.

\section*{7. Immediate Next Step}

The next valid step is:

\begin{itemize}
\item run the full locked 3-seed scaled Push-T suite,
\item aggregate the five requested configuration results,
\item issue a continue/stop recommendation from those scaled numbers,
\item add the feasibility critic only if the hierarchy advantage survives that scaled comparison.
\end{itemize}

\section*{8. Bottom Line}

This round completed the locked reevaluation and the missing plumbing needed for the scaled suite.
It did \textbf{not} complete the scaled suite itself.
The current best debug checkpoint still shows hierarchical value, but only as a distance and latency advantage.
The line should remain frozen until the scaled pilot produces the actual continue/stop evidence.

\end{document}
